\documentclass[12pt]{article}

\usepackage[a4paper, margin=2.5cm]{geometry}
\usepackage{amsmath}
\usepackage{amssymb}
\usepackage{amsthm}
\usepackage[colorlinks=true, linkcolor=blue, citecolor=blue, urlcolor=blue]{hyperref}
\usepackage{booktabs}
\usepackage{natbib}

\newtheorem{theorem}{Theorem}
\newtheorem{proposition}[theorem]{Proposition}
\newtheorem{corollary}[theorem]{Corollary}
\newtheorem{definition}[theorem]{Definition}
\newtheorem{remark}[theorem]{Remark}
\newtheorem{example}[theorem]{Example}
\newtheorem{lemma}[theorem]{Lemma}

\newcommand{\Ecal}{\mathcal{E}}
\newcommand{\Ocal}{\mathcal{O}}
\newcommand{\Gcal}{\mathcal{G}}

\title{Metric Emergence from Event-Order Geometry}
\author{%
  Lee Smart\\[4pt]
  \textit{Institute of Vibrational Field Dynamics}\\[2pt]
  \texttt{contact@vibrationalfielddynamics.org}\\[2pt]
  \texttt{@vfd\_org}%
}
\date{April 2026}

\begin{document}

\maketitle

\noindent\textbf{Status:} Preprint (not peer-reviewed)

\begin{abstract}
Papers~XXIII--XXIV derived a strict partial order on admissible overlap events and showed that observer-dependent time functions, relativity of simultaneity, and time-dilation analogues emerge from the event poset. The present paper addresses the next question: what does it mean for two events to be \textit{near or far}?

Three candidate separation constructions are developed and compared. For comparable events (those related by the partial order), a \textit{cover-step (chain-length) separation} counts the minimum number of cover steps along a saturated chain between them, providing an observer-independent directed temporal separation. For incomparable events (those not causally related), a \textit{family-relative observer-disagreement separation} measures the variance of time-difference assignments across a chosen finite observer family, providing a structural precursor to spatial distance relative to that family. A \textit{transition-cost separation} on the Hasse graph gives a unified graph-distance that applies to both cases on a connected component of the graph.

The cover-step separation is defined in terms of saturated chains and its basic order-theoretic properties (non-negativity, zero-on-diagonal, and a forward triangle inequality on comparable triples) are proved. The observer-disagreement construction is shown to have sign of the individual time-differences that is family-independent for comparable pairs, and possibly sign-varying for incomparable pairs; the resulting variance is strictly positive whenever the chosen family contains sign-reversing observers for the pair. On each connected component of the Hasse graph, the transition-cost construction satisfies all metric axioms. A worked example on the first two birth-angle layers of the dual-pentagon system of Papers~XXIII--XXIV demonstrates all three constructions explicitly.

The paper does not derive the Lorentz metric or continuous geometry. It establishes that the event poset canonically yields a graph metric on each connected Hasse-graph component, and provides two additional, weaker separation functions (cover-step and family-relative observer-disagreement) that together distinguish comparable from incomparable event pairs.
\end{abstract}

%==================================================================
\section{Introduction}\label{sec:intro}
%==================================================================

Paper~XXIII~\citep{SmartXXIII} derived a strict partial order $\prec$ on an event set $\Ecal$ from phase-overlap geometry. Paper~XXIV~\citep{SmartXXIV} showed that observers (linear extensions with sampling measures) produce frame-dependent time functions, relativity of simultaneity, and time-dilation analogues.

\medskip
\noindent\textbf{Standing assumption.} Throughout this paper, $\Ecal$ denotes a finite event set carrying a strict partial order $\prec$ (as in Papers~XXIII--XXIV). Finiteness guarantees that every comparable pair $e_1 \prec e_2$ admits at least one saturated chain from $e_1$ to $e_2$ (Definition~\ref{def:satchain}) and that the maximum chain length between such a pair is finite.
\medskip


These results give the event set temporal and observer structure, but not spatial or metric structure. The present paper asks: \textit{given only the event poset and its observer completions, can one define a meaningful notion of separation between events?}

\subsection*{The essential distinction}

Two events can be related in exactly two ways under $\prec$:
\begin{itemize}
    \item \textbf{Comparable} ($e_1 \prec e_2$ or $e_2 \prec e_1$): causally / temporally ordered.
    \item \textbf{Incomparable} (neither holds): ``simultaneous'' / spacelike-separated.
\end{itemize}

Any separation construction must handle both cases, and the two cases may require different structures. The paper develops three candidates and compares them.

\subsection*{Scope}

This paper constructs separation quantities and proves their basic properties. It does not derive the Lorentz metric, continuous manifold structure, or the Einstein equations. It establishes a graph metric on each connected Hasse-graph component, together with two weaker separation constructions (cover-step and family-relative observer-disagreement).

%==================================================================
\section{Cover Relations and the Transition Graph}\label{sec:hasse}
%==================================================================

\begin{definition}[Cover relation]\label{def:cover}
Event $e_1$ is \textit{covered by} $e_2$ (written $e_1 \lessdot e_2$) if $e_1 \prec e_2$ and there is no event $e$ with $e_1 \prec e \prec e_2$.
\end{definition}

\begin{definition}[Transition graph (Hasse graph)]\label{def:hasse}
The \textit{transition graph} $\Gcal = (\Ecal, E_H)$ has vertex set $\Ecal$ and edge set
\begin{equation}
    E_H = \{(e_1, e_2) : e_1 \lessdot e_2\}
\end{equation}
This is the Hasse diagram of $(\Ecal, \prec)$, viewed as an undirected graph (ignoring edge orientation for the purpose of graph distance).
\end{definition}

\begin{remark}
The cover relation encodes ``immediate succession'': $e_1 \lessdot e_2$ means $e_2$ is the next event after $e_1$ with no intervening event in the partial order. The transition graph captures the minimal-step accessibility structure of the event set.
\end{remark}

%==================================================================
\section{Construction A: Chain-Length Separation}\label{sec:chain}
%==================================================================

For comparable events, the natural separation is the number of order steps between them.

\begin{definition}[Saturated chain]\label{def:satchain}
A \textit{saturated chain} from $e_1$ to $e_2$ is a sequence $e_1 = c_0, c_1, \ldots, c_k = e_2$ in $\Ecal$ with $c_{i-1} \lessdot c_i$ for all $i = 1, \ldots, k$: i.e.\ every consecutive pair is a cover pair, with no intermediate events. Its \textit{length} is $k$.
\end{definition}

\begin{definition}[Cover-step (chain-length) separation]\label{def:dchain}
The \textit{natural domain} of $d_{\mathrm{chain}}$ is $\{(e,e) : e \in \Ecal\} \cup \{(e_1, e_2) \in \Ecal \times \Ecal : e_1 \prec e_2\}$ --- the diagonal together with strictly comparable pairs. On this domain, define:
\begin{equation}\label{eq:dchain}
    d_{\mathrm{chain}}(e, e) = 0, \qquad d_{\mathrm{chain}}(e_1, e_2) = \min\{k \in \mathbb{Z}_{\geq 1} : \text{a saturated chain of length } k \text{ from } e_1 \text{ to } e_2 \text{ exists}\}.
\end{equation}
Since $\Ecal$ is finite (standing assumption), every strictly comparable pair admits at least one saturated chain, so $d_{\mathrm{chain}}$ is defined everywhere on its natural domain and takes values in $\mathbb{Z}_{\geq 0}$.
\end{definition}

\begin{remark}
$d_{\mathrm{chain}}(e_1, e_2) = k$ means there is a saturated chain $e_1 = c_0 \lessdot c_1 \lessdot \cdots \lessdot c_k = e_2$ of exactly $k$ cover steps, and no saturated chain with fewer cover steps exists. Equivalently, $d_{\mathrm{chain}}$ is the \textit{Hasse-graph distance restricted to directed (upward) paths}. An arbitrary chain $e_1 \prec c \prec e_2$ in which consecutive elements are not cover-related does \emph{not} decrease $d_{\mathrm{chain}}$, because such a chain can always be refined (in finite $\Ecal$) to a longer saturated chain.
\end{remark}

\begin{proposition}[Properties of chain-length separation]\label{prop:dchain}
On its natural domain (Definition~\ref{def:dchain}), $d_{\mathrm{chain}}$ satisfies:
\begin{enumerate}
    \item \textbf{Non-negativity:} $d_{\mathrm{chain}}(e_1, e_2) \geq 0$.
    \item \textbf{Identity:} $d_{\mathrm{chain}}(e, e) = 0$.
    \item \textbf{Triangle inequality:} If $e_1 \prec e_2 \prec e_3$:
    \begin{equation}
        d_{\mathrm{chain}}(e_1, e_3) \leq d_{\mathrm{chain}}(e_1, e_2) + d_{\mathrm{chain}}(e_2, e_3)
    \end{equation}
\end{enumerate}
\end{proposition}

\begin{proof}
(1): Saturated-chain lengths are non-negative integers. (2): The trivial saturated chain consisting of the single event $e$ has length $0$. (3): Concatenating a shortest saturated chain from $e_1$ to $e_2$ with one from $e_2$ to $e_3$ yields a saturated chain from $e_1$ to $e_3$ (the last cover step of the first chain ends at $e_2$, and the first cover step of the second begins at $e_2$) of length $d_{\mathrm{chain}}(e_1, e_2) + d_{\mathrm{chain}}(e_2, e_3)$. The minimum over all saturated chains from $e_1$ to $e_3$ can only be smaller or equal.
\end{proof}

\begin{remark}
$d_{\mathrm{chain}}$ is a directed separation function on its natural domain (diagonal plus strictly comparable pairs), not a metric: it is asymmetric --- $d_{\mathrm{chain}}(e_2, e_1)$ is undefined whenever $e_1 \prec e_2$ --- and it is not defined on incomparable pairs at all. This separation is \textit{temporal}: it measures causal distance, not spatial distance.
\end{remark}

\begin{remark}[Longest chain as proper-time analogue]
In causal set theory~\citep{BombelliLeeMeyerSorkin1987,BrightwellGregory1991}, the longest (not shortest) chain between comparable events is the analogue of proper time. Define:
\begin{equation}
    d_{\max}(e_1, e_2) = \max\{k : \text{there exists a saturated chain of length } k \text{ from } e_1 \text{ to } e_2\}
\end{equation}
$d_{\max}$ is well-defined and finite for comparable pairs in a finite $\Ecal$, but satisfies a \emph{reverse} triangle inequality ($d_{\max}(e_1, e_3) \geq d_{\max}(e_1, e_2) + d_{\max}(e_2, e_3)$ when $e_1 \prec e_2 \prec e_3$) rather than a metric-style triangle inequality. It is better viewed as a temporal-size functional than a metric. In the VFD framework, whether $d_{\mathrm{chain}}$ (shortest saturated) or $d_{\max}$ (longest saturated) better approximates physical proper time is an open question; in causal set theory, the longest chain is preferred.
\end{remark}

%==================================================================
\section{Construction B: Observer-Disagreement Separation}\label{sec:disagreement}
%==================================================================

For incomparable events, no chain exists and $d_{\mathrm{chain}}$ is undefined. A different construction is needed.

\begin{definition}[Family-relative observer-disagreement separation]\label{def:dobs}
Fix a finite observer family $\Ocal = \{O_1, \ldots, O_N\}$ with $N \geq 2$, each observer $O_i$ as in Paper~XXIV (a linear extension of $\prec$ together with a positive weight function on $\Ecal$). Endow $\Ocal$ with the uniform (counting) measure so that expectations and variances are population statistics. For any $e_1, e_2 \in \Ecal$, define
\begin{equation}\label{eq:dobs}
    \sigma_{\Ocal}(e_1, e_2) \;=\; \frac{1}{N}\sum_{i=1}^{N} \bigl(\Delta_i - \bar{\Delta}\bigr)^2, \qquad \Delta_i := t_{O_i}(e_1) - t_{O_i}(e_2), \quad \bar{\Delta} := \frac{1}{N}\sum_{i=1}^{N}\Delta_i.
\end{equation}
$\sigma_{\Ocal}$ depends on the chosen family $\Ocal$; the same event pair can yield different values under different families.
\end{definition}

\begin{proposition}[Family-relative sign structure]\label{prop:dobs}
Let $\Ocal$ be a finite observer family as in Definition~\ref{def:dobs}, and let $\Delta_i$ be as in~\eqref{eq:dobs}.
\begin{enumerate}
    \item If $e_1 = e_2$: $\Delta_i = 0$ for every $O_i \in \Ocal$; hence $\sigma_{\Ocal}(e_1, e_2) = 0$.
    \item If $e_1 \prec e_2$ (comparable): $\Delta_i < 0$ for every $O_i \in \Ocal$; the sign of $\Delta_i$ is family-independent.
    \item If $e_1$ and $e_2$ are incomparable, there exist linear extensions placing them in either relative order (Paper~XXIV, relativity-of-simultaneity theorem, \citealt{SmartXXIV}). In particular, if $\Ocal$ is chosen to contain at least one observer with $\Delta_i < 0$ and at least one with $\Delta_j > 0$, then the $\Delta_i$ take both signs across $\Ocal$, which forces $\sigma_{\Ocal}(e_1, e_2) > 0$.
\end{enumerate}
\end{proposition}

\begin{proof}
(1): Immediate from $t_O(e) = t_O(e)$. (2): $e_1 \prec e_2$ forces $t_O(e_1) < t_O(e_2)$ for every observer by Paper~XXIV's order-preservation proposition~\citep{SmartXXIV}. (3): If the $\{\Delta_i\}_{i=1}^N$ take strictly positive and strictly negative values, they are not all equal, so the population variance of a non-constant real sample is strictly positive.
\end{proof}

\begin{remark}[Interpretation]
The family-relative observer-disagreement $\sigma_{\Ocal}$ serves as a \textit{spatial-separation precursor relative to the chosen family}: it is large when observers in $\Ocal$ strongly disagree on the time-difference, and zero whenever all observers in $\Ocal$ assign the same $\Delta$ (including, but not limited to, the case $e_1 = e_2$). Comparable events have a family-independent sign of $\Delta_i$; for incomparable events, $\sigma_{\Ocal}$ can be made strictly positive by choosing $\Ocal$ to include a sign-reversing pair, but it is not a function of the event pair alone.
\end{remark}

\begin{remark}[Not a metric]
$\sigma_{\Ocal}$ does not satisfy the triangle inequality in general and is therefore not a metric. It is a meaningful \textit{separation measure relative to a chosen observer family}, not a distance function in the strict sense. Its role is to provide nontrivial observer-sensitive structure for incomparable pairs, where $d_{\mathrm{chain}}$ is undefined.
\end{remark}

%==================================================================
\section{Construction C: Transition-Cost Metric}\label{sec:transition}
%==================================================================

The transition graph $\Gcal$ (Definition~\ref{def:hasse}) provides a unified metric that applies to both comparable and incomparable events.

\begin{definition}[Transition-cost metric]\label{def:dtrans}
For $e_1, e_2 \in \Ecal$ in the same connected component of $\Gcal$, define
\begin{equation}\label{eq:dtrans}
    d_{\mathrm{trans}}(e_1, e_2) = \text{shortest-path length in } \Gcal \text{ between } e_1 \text{ and } e_2,
\end{equation}
where $\Gcal$ is the undirected Hasse graph with unit edge weights. If $\Gcal$ is connected (a hypothesis verified in the pentagon model below but not established in general for the full 600-cell event set), $d_{\mathrm{trans}}$ is defined on all of $\Ecal \times \Ecal$.
\end{definition}

\begin{theorem}[Metric on a connected component]\label{thm:metric}
On any connected component of $\Gcal$, $d_{\mathrm{trans}}$ is a metric:
\begin{enumerate}
    \item $d_{\mathrm{trans}}(e_1, e_2) \geq 0$, with equality iff $e_1 = e_2$.
    \item $d_{\mathrm{trans}}(e_1, e_2) = d_{\mathrm{trans}}(e_2, e_1)$.
    \item $d_{\mathrm{trans}}(e_1, e_3) \leq d_{\mathrm{trans}}(e_1, e_2) + d_{\mathrm{trans}}(e_2, e_3)$.
\end{enumerate}
If $\Gcal$ is connected, $d_{\mathrm{trans}}$ is a metric on the full event set $\Ecal$.
\end{theorem}

\begin{proof}
These are the standard properties of shortest-path distance on a connected graph. (1): Shortest paths have non-negative length; length $0$ iff $e_1 = e_2$. (2): Undirected graph, so paths are reversible. (3): Concatenation of shortest paths from $e_1$ to $e_2$ and from $e_2$ to $e_3$ gives an upper bound on the $e_1$-to-$e_3$ path length.
\end{proof}

\begin{remark}
$d_{\mathrm{trans}}$ is the first genuine metric on (each connected component of) the event set. Unlike $d_{\mathrm{chain}}$ (which requires comparability) and $\sigma_{\Ocal}$ (which fails the triangle inequality), $d_{\mathrm{trans}}$ applies to all same-component pairs and satisfies all metric axioms.
\end{remark}

\begin{remark}[Comparable vs incomparable]
For comparable events $e_1 \prec e_2$, every saturated chain from $e_1$ to $e_2$ gives a path in $\Gcal$ of the same length (each cover step is a Hasse edge). Therefore
\begin{equation}
    d_{\mathrm{trans}}(e_1, e_2) \;\leq\; d_{\mathrm{chain}}(e_1, e_2).
\end{equation}
In the pentagon toy model (Section~\ref{sec:example}), equality holds: $d_{\mathrm{trans}} = d_{\mathrm{chain}} = 1$ for all comparable pairs. Whether equality holds in general depends on the specific Hasse-graph structure. For incomparable events, $d_{\mathrm{trans}}$ involves ``sideways'' steps through the graph --- paths that go up then down (or vice versa) through the partial order. These sideways paths are the structural precursor to spatial separation.
\end{remark}

%==================================================================
\section{Temporal vs Spatial Separation}\label{sec:split}
%==================================================================

The three constructions cleanly separate into temporal and spatial-precursor components:

\begin{table}[ht]
\centering
\caption{Separation structure by event relationship.}
\label{tab:separation}
\begin{tabular}{@{}llll@{}}
\toprule
Event pair & $d_{\mathrm{chain}}$ & $\sigma_{\Ocal}$ (family-relative) & $d_{\mathrm{trans}}$ \\
\midrule
Comparable & Defined (cover steps) & Sign of $\Delta$ family-independent & Graph distance $\leq d_{\mathrm{chain}}$ \\
Incomparable & Undefined & Sign of $\Delta$ can vary; $\sigma_{\Ocal}>0$ for sign-reversing $\Ocal$ & Sideways path \\
\midrule
Interpretation & Temporal interval & Family-relative spatial precursor & Unified metric (on component) \\
\bottomrule
\end{tabular}
\end{table}

\begin{remark}[Analogy to Lorentzian structure]
In Lorentzian geometry, timelike intervals are computed along causal curves, and spacelike intervals require leaving the light cone. The VFD framework exhibits a structural analogue: comparable events have chain-based (causal-curve) separation, while incomparable events require sideways (off-chain) paths in the transition graph. The transition-cost metric $d_{\mathrm{trans}}$ does not distinguish timelike from spacelike --- it is a single positive-definite distance. A signed interval $s^2 = d_{\mathrm{chain}}^2 - d_{\mathrm{sideways}}^2$ could in principle be defined but is not pursued here.
\end{remark}

%==================================================================
\section{Observer Time vs Structural Interval}\label{sec:comparison}
%==================================================================

\begin{proposition}[Structural bound on observer time]\label{prop:bound}
For comparable events $e_1 \prec e_2$ and any observer $O = (\leq_O, \mu_O)$:
\begin{equation}
    t_O(e_2) - t_O(e_1) \geq d_{\mathrm{chain}}(e_1, e_2) \cdot \min_{e \in \Ecal} \mu_O(e)
\end{equation}
\end{proposition}

\begin{proof}
Let $e_1 = c_0 \lessdot c_1 \lessdot \cdots \lessdot c_k = e_2$ be a shortest saturated chain from $e_1$ to $e_2$ (so $k = d_{\mathrm{chain}}(e_1, e_2)$; Definition~\ref{def:dchain}). Since $\leq_O$ extends $\prec$ and each cover step gives $c_{i-1} \prec c_i$, the observer ordering satisfies $c_0 <_O c_1 <_O \cdots <_O c_k$. Therefore the $k$ events $c_1, \ldots, c_k$ lie in the interval $\{e' \in \Ecal : e_1 <_O e' \leq_O e_2\}$, which is the set $S$ in Paper~XXIV's observer-endpoint-interval construction~\citep{SmartXXIV}; each of them contributes at least $\min_{e \in \Ecal} \mu_O(e) > 0$ to the sum for $t_O(e_2) - t_O(e_1)$.
\end{proof}

\begin{remark}
This proposition connects observer-dependent time to observer-independent structure: the chain-length separation provides a lower bound on any observer's elapsed time, scaled by the observer's minimum sampling rate. The structural interval ``lives underneath'' all observer times.
\end{remark}

%==================================================================
\section{Worked Example: Dual Pentagons}\label{sec:example}
%==================================================================

\begin{example}[Pentagon separation structure]\label{ex:pentagon}
Using the first two birth-angle layers of the dual-pentagon event set of Paper~XXIII~\citep{SmartXXIII} (also used in Paper~XXIV~\citep{SmartXXIV}):
\[
    L_1 = \{(0,4),(1,0),(2,1),(3,2),(4,3)\} \quad \text{at } \theta_b = \pi/5 - \delta,
\]
\[
    L_2 = \{(0,3),(1,4),(2,0),(3,1),(4,2)\} \quad \text{at } \theta_b = 3\pi/5 - \delta,
\]
ten events in two layers.

\textbf{Partial order:} Every $L_1$ event precedes every $L_2$ event (layer-1 birth angles are strictly smaller); events within a single layer share a birth angle and are incomparable.

\textbf{Cover relations:} Within the first-two-layer truncation of the event set, every $L_2$ event covers every $L_1$ event, i.e.\ $e \lessdot f$ for every $e \in L_1$ and $f \in L_2$ (Definition~\ref{def:cover}; no events lie strictly between the two layers in this truncation).

\textbf{Hasse graph:} The complete bipartite graph $K_{5,5}$ between $L_1$ and $L_2$; it is connected, so Theorem~\ref{thm:metric} applies on the full 10-event set.

\subsection*{Chain-length separation (comparable pairs)}

For any $e \in L_1$ and $f \in L_2$: there is a one-step saturated chain $e \lessdot f$, so $d_{\mathrm{chain}}(e, f) = 1$. All comparable pairs have the same chain-length separation.

\subsection*{Family-relative observer-disagreement separation (incomparable pairs)}

Within $L_1$, the events $(0,4)$ and $(1,0)$ are incomparable. Take the two observers from Paper~XXIV's Example (same event set, uniform weight $\mu \equiv 1$): $O_1$ uses the linear extension $\leq_A$, $O_2$ uses $\leq_B$ obtained from $\leq_A$ by swapping the incomparable pair $(0,4)$ and $(1,0)$. Form the family $\Ocal = \{O_1, O_2\}$. Then
\begin{itemize}
    \item $O_1$: $t_{O_1}(0,4) = 1$, $t_{O_1}(1,0) = 2$, so $\Delta_1 = t_{O_1}(0,4) - t_{O_1}(1,0) = -1$.
    \item $O_2$: $t_{O_2}(0,4) = 2$, $t_{O_2}(1,0) = 1$, so $\Delta_2 = t_{O_2}(0,4) - t_{O_2}(1,0) = +1$.
\end{itemize}
Mean $\bar\Delta = 0$, variance $\sigma_{\Ocal}((0,4),(1,0)) = \tfrac{1}{2}\bigl((-1)^2 + (+1)^2\bigr) = 1$. The incomparable pair is separated by the family $\Ocal$ in the sense of Proposition~\ref{prop:dobs}(3); the value itself depends on the family. By contrast, for any comparable $e \in L_1$, $f \in L_2$ and the same $\Ocal$, the individual $\Delta_i$ are all strictly negative, so the sign of $\Delta$ is family-independent (Proposition~\ref{prop:dobs}(2)).

\subsection*{Transition-cost metric (all pairs)}

The Hasse graph is $K_{5,5}$, so:
\begin{itemize}
    \item Comparable pair (different layers): $d_{\mathrm{trans}} = 1$ (direct Hasse edge).
    \item Incomparable pair (same layer): $d_{\mathrm{trans}} = 2$. For example, $(0,4) \to (0,3) \to (1,0)$ is a length-2 path in $K_{5,5}$, via any common neighbour.
\end{itemize}

\subsection*{Summary}

\begin{table}[ht]
\centering
\caption{Separations in the first-two-layer pentagon truncation, with $\Ocal = \{O_1, O_2\}$ as above.}
\label{tab:pentagon}
\begin{tabular}{@{}llll@{}}
\toprule
Pair type & $d_{\mathrm{chain}}$ & $\sigma_{\Ocal}$ & $d_{\mathrm{trans}}$ \\
\midrule
Same event & $0$ & $0$ & $0$ \\
Comparable ($L_1 \to L_2$) & $1$ & $\Delta_i < 0$ for all $i$ (fixed sign) & $1$ \\
Incomparable (same layer) & undefined & $1$ (for the incomparable $(0,4),(1,0)$) & $2$ \\
\bottomrule
\end{tabular}
\end{table}

In this corrected pentagon example, the transition-cost metric gives $d_{\mathrm{trans}} = 1$ for ``timelike'' (comparable) pairs and $d_{\mathrm{trans}} = 2$ for ``spacelike'' (incomparable) pairs. Incomparable events are \textit{farther} in the transition metric than comparable ones --- the sideways cost exceeds the forward cost.
\end{example}

%==================================================================
\section{Relation to Known Frameworks}\label{sec:comparison_lit}
%==================================================================

\begin{table}[ht]
\centering
\caption{Comparison with causal-set and graph-metric approaches.}
\label{tab:comparison}
\begin{tabular}{@{}lll@{}}
\toprule
Feature & VFD (this paper) & Causal set / graph theory \\
\midrule
Temporal separation & Min-chain length & Longest chain (proper time) \\
Spatial precursor & Observer disagreement & Volume-based estimators \\
Unified metric & Hasse graph distance & Graph metric on Hasse diagram \\
Metric derivation & From event-order geometry & From assumed causal structure \\
Signed interval & Not yet constructed & Lorentzian distance (assumed) \\
\bottomrule
\end{tabular}
\end{table}

\begin{remark}
In causal set theory, the longest chain between comparable events approximates Lorentzian proper time, and spatial distance is estimated from causal-diamond volumes. The VFD approach differs in that the event set and order are \textit{derived} from 600-cell geometry rather than postulated. The metric constructions here are structural parallels, not claimed equivalences.
\end{remark}

%==================================================================
\section{Discussion and Limitations}\label{sec:discussion}
%==================================================================

\subsection{What is established}

\begin{itemize}
    \item Cover relation, saturated chain, and Hasse graph explicitly defined (Definitions~\ref{def:cover}, \ref{def:hasse}, \ref{def:satchain}).
    \item Cover-step (chain-length) separation $d_{\mathrm{chain}}$: well-defined directed temporal separation on comparable pairs in finite $\Ecal$ (Proposition~\ref{prop:dchain}).
    \item Family-relative observer-sign structure (Proposition~\ref{prop:dobs}): for comparable pairs, the sign of $\Delta_i$ is family-independent; for incomparable pairs, the sign can be made to vary across a chosen family, yielding $\sigma_{\Ocal} > 0$. The variance $\sigma_{\Ocal}$ depends on the chosen family $\Ocal$ and does not satisfy the triangle inequality.
    \item Transition-cost metric: a genuine metric on each connected component of the Hasse graph $\Gcal$ (Theorem~\ref{thm:metric}); a metric on the full $\Ecal$ only if $\Gcal$ is connected.
    \item Structural bound relating observer elapsed time to $d_{\mathrm{chain}}$ (Proposition~\ref{prop:bound}).
    \item All constructions demonstrated on the first-two-layer truncation of the XXIII/XXIV pentagon toy model.
\end{itemize}

\subsection{What is not established}

\begin{itemize}
    \item No Lorentzian metric. The transition-cost metric is positive-definite, not Lorentzian.
    \item No signed interval. A $ds^2 = dt^2 - dx^2$ decomposition is not derived.
    \item No continuous manifold. All constructions are discrete (graph-based).
    \item No Einstein equations or curvature. No curvature notion is introduced in this paper, even for the toy model.
    \item The observer-disagreement quantity $\sigma_{\Ocal}$ is family-dependent and not a metric (triangle inequality failure is explicit).
    \item The Hasse graph connectivity required for a global metric on all of $\Ecal$ (Theorem~\ref{thm:metric}) is verified in the pentagon truncation but is not established for the full 600-cell event structure.
\end{itemize}

\subsection{Metric failure modes (honest assessment)}

\begin{itemize}
    \item $d_{\mathrm{chain}}$: only defined for comparable pairs; it is a directed temporal separation built from cover steps, not a metric on full $\Ecal$.
    \item $\sigma_{\Ocal}$: depends on the chosen observer family $\Ocal$, and fails the triangle inequality; not a metric in the formal sense.
    \item $d_{\mathrm{trans}}$: satisfies all metric axioms on any connected component of $\Gcal$ but is positive-definite, losing the timelike/spacelike distinction that Lorentzian geometry requires.
\end{itemize}

A future construction that recovers a signed interval from the interplay of $d_{\mathrm{chain}}$ and $d_{\mathrm{trans}}$ would represent a significant step toward Lorentzian structure. The present paper provides the ingredients.

%==================================================================
\section{Conclusion}\label{sec:conclusion}
%==================================================================

The event-order framework of Papers~XXIII--XXIV supports emergent separation structure. Three constructions are developed:
\begin{itemize}
    \item \textbf{Cover-step (chain-length) separation} $d_{\mathrm{chain}}$: an observer-independent directed temporal separation for comparable events, defined as the minimum number of cover steps along a saturated chain (properties proved).
    \item \textbf{Family-relative observer-disagreement separation} $\sigma_{\Ocal}$: a nontrivial separation measure for incomparable events relative to a chosen finite observer family, interpreted as a family-relative spatial precursor (sign structure proved; not a metric; depends on the chosen family).
    \item \textbf{Transition-cost metric} $d_{\mathrm{trans}}$: a metric on each connected component of the Hasse graph, obtained from shortest-path distance (metric axioms proved; global on $\Ecal$ only if $\Gcal$ is connected).
\end{itemize}

The structural distinction between comparable (chain-based, temporal) and incomparable (sideways, spatial-precursor) separation provides a discrete structural analogue of the comparable/incomparable split familiar from Lorentzian geometry, though a signed interval is not yet derived.

The event-order framework now carries: temporal ordering (Paper~XXIII), observer-dependent kinematics (Paper~XXIV), and separation geometry (this paper). What remains for the GR programme is the construction of curvature and dynamical equations on this structure.

%==================================================================
\bibliographystyle{plainnat}
\bibliography{references}

\end{document}
